\chapter{HTLV-1 Infection}

\section*{Definition and Overview}
Human T-lymphotropic virus type 1 (HTLV-1) is a retrovirus that infects T-lymphocytes, a type of white blood cell. It is related to HIV but behaves differently: most people infected with HTLV-1 remain asymptomatic for life. In a minority, however, the virus causes serious diseases, including adult T-cell leukaemia/lymphoma (ATLL) and HTLV-1-associated myelopathy (HAM), a progressive spinal cord disease. HTLV-1 is transmitted through blood, sexual contact, breast milk, and contaminated needles. There is no cure or vaccine, but most infected people never develop disease.

\section*{Epidemiology}
HTLV-1 infects an estimated 5--10 million people worldwide, with high prevalence in parts of Japan, the Caribbean, Central and South America, West and Central Africa, the Middle East, and in Indigenous communities in central Australia. In Australia, prevalence is highest among Aboriginal and Torres Strait Islander people in remote central Australia, where it is a significant health issue. About 5--10\% of infected people develop ATLL or HAM over their lifetime, usually decades after infection. Most carriers remain well.

\section*{Aetiology and Pathophysiology}
\begin{itemize}
    \item \textbf{Transmission:} breastfeeding (the main route in children), unprotected sex, blood transfusion with infected blood, and sharing contaminated injecting equipment
    \item \textbf{Pathophysiology:} the virus infects and immortalises CD4 T-cells, integrating into the host genome
    \item \textbf{Adult T-cell leukaemia/lymphoma:} malignant transformation of infected T-cells causes an aggressive blood cancer
    \item \textbf{HTLV-1-associated myelopathy:} immune and viral mechanisms damage the spinal cord, causing progressive leg weakness
    \item \textbf{Other effects:} uveitis, dermatitis, and increased susceptibility to some infections
    \item \textbf{Most carriers:} remain asymptomatic throughout life
    \item \textbf{Incubation:} disease develops decades after infection
\end{itemize}

\section*{Clinical Features}
\begin{itemize}
    \item Most people have no symptoms and are diagnosed on blood tests
    \item ATLL: enlarged lymph nodes, fever, skin lesions, and blood abnormalities
    \item HAM: progressive stiffness and weakness of the legs, back pain, and bladder problems
    \item Walking difficulty developing over months to years
    \item Uveitis: eye pain, redness, and blurred vision
    \item Recurrent infections due to impaired immunity
    \item Symptoms vary with the disease the virus causes
\end{itemize}

\section*{Differential Diagnosis}
\begin{itemize}
    \item Other causes of myelopathy, including multiple sclerosis, vitamin B12 deficiency, and spinal compression
    \item Other lymphomas and leukaemias in ATLL
    \item Other retroviruses, including HIV
    \item The diagnosis is confirmed by blood testing for HTLV-1 antibodies and viral DNA
    \item Most people with positive tests have no disease and need monitoring
\end{itemize}

\section*{When to Seek Medical Care}
Anyone at risk should consider testing, including people from high-prevalence regions, partners of infected people, and those who received blood products before screening began. Symptoms such as progressive leg weakness, eye problems, or enlarged nodes need review.

\textbf{Call 000 (or local emergency services) for:}
\begin{itemize}
    \item Sudden inability to walk or severe weakness
    \item High fever with enlarged lymph nodes and feeling very unwell
    \item Difficulty breathing or chest pain
    \item Collapse or confusion
    \item Any life-threatening emergency
\end{itemize}

\section*{Diagnosis and Investigations}
\begin{itemize}
    \item \textbf{Blood tests:} HTLV-1 antibody testing, confirmed with viral DNA tests
    \item \textbf{Blood screening:} blood donations are screened for HTLV-1 in Australia
    \item \textbf{Full blood count and cell markers:} for suspected ATLL
    \item \textbf{MRI of the spine:} for suspected HAM
    \item \textbf{Lumbar puncture:} cerebrospinal fluid analysis in neurological disease
    \item \textbf{Monitoring:} regular review of carriers for early signs of disease
\end{itemize}

\section*{Management}
\begin{itemize}
    \item \textbf{Carrier monitoring:} regular clinical review and blood tests
    \item \textbf{Prevention of transmission:} avoid breastfeeding where safe alternatives exist, use condoms, and do not share needles
    \item \textbf{ATLL treatment:} chemotherapy and newer targeted therapies, with specialist haematology care
    \item \textbf{HAM treatment:} symptom management, physiotherapy, and medicines to slow progression
    \item \textbf{Complications:} management of infections and eye disease
    \item \textbf{Counselling:} information and support for carriers and families
    \item \textbf{Research:} clinical trials of treatments and vaccines are ongoing
\end{itemize}

\section*{Complications}
\begin{itemize}
    \item Adult T-cell leukaemia/lymphoma, which is often aggressive
    \item HTLV-1-associated myelopathy with progressive disability
    \item Uveitis and other inflammatory conditions
    \item Increased risk of certain infections
    \item Transmission to partners and children
    \item Psychological impact of a chronic infection
\end{itemize}

\section*{Prognosis and Outcomes}
The majority of people with HTLV-1 never develop disease and live normal lives. For the 5--10\% who do, outcomes depend on the condition: ATLL is serious and requires intensive treatment, while HAM progresses gradually and is managed symptomatically. Research into treatments is advancing, and specialist care improves outcomes. Avoiding transmission protects future generations, particularly through the prevention of breastfeeding transmission in affected communities.

\section*{Prevention and Self-Care}
\begin{itemize}
    \item Discuss testing with your doctor if you are from a high-prevalence region or have risk factors
    \item Use condoms to prevent sexual transmission
    \item Never share needles or injecting equipment
    \item Infected mothers should discuss infant feeding options
    \item Blood donors are screened, so blood products are safe
    \item Attend regular monitoring if you are a carrier
    \item Report new neurological or blood symptoms promptly
    \item Seek specialist care at a centre experienced with HTLV-1
\end{itemize}

\section*{Sources and Further Reading}
The following reputable sources provide further information for healthcare students and patients seeking reliable, current advice about HTLV-1 infection.

\begin{itemize}
    \item Healthdirect. \textit{HTLV-1 infection.} https://www.healthdirect.gov.au/htlv
    \item ASHM. \textit{HTLV-1 testing and clinical guidance.} https://www.ashm.org.au/
    \item Wardliparingga Aboriginal Health Equity and other research centres. \textit{HTLV-1 in central Australia.} https://www.sahmri.org.au/
    \item NSW Health. \textit{HTLV-1 fact sheet.} https://www.health.nsw.gov.au/
    \item World Health Organization. \textit{HTLV-1 information.} https://www.who.int/
    \item International Retrovirology Association. \textit{HTLV research and education.} https://htlv.org/
\end{itemize}
